\begin{ExerciseBlock}
\ExerciseHeading{10}{Beschleunigter Bahnwagen}{forces_accelerating_railcar}
In einem Eisenbahnwagen sei eine Bleikugel mit Masse $m$ an einem langen Faden aufgehängt. Welchen Winkel $\varphi$ wird der Faden mit der Vertikalen in der Gleichgewichtslage des Pendels für die folgenden 3 Fälle einschließen?

\begin{enumerate}
	\item Während des Anfahrens: Die Beschleunigung sei geradlinig gleichförmig mit einem Betrag, der nach Durchlaufen von $2\,$km aus dem Zustand der Ruhe eine Geschwindigkeit $v_0 = 72\,$km/h liefert.

	\item Während der Fahrt auf gerader Strecke mit konstanter Geschwindigkeit $v_0 = 72\,$km/h.

	\item Während des Durchlaufens einer Kurve mit konstantem Betrag der Geschwindigkeit $v_0 = 72\,$km/h. Der Krümmungsradius der Kurve sei $r_0 = 1\,$km.

\end{enumerate}
\end{ExerciseBlock}
\UseExerciseSolution{forces_accelerating_railcar}
